\documentclass[11pt,a4paper]{article}
\usepackage[utf8]{inputenc}
\usepackage{amsmath,amssymb,amsfonts}
\usepackage{graphicx}
\usepackage{geometry}
\usepackage{enumitem}
\usepackage{fancyhdr}

\geometry{top=1in, bottom=1in, left=0.8in, right=0.8in}
\pagestyle{fancy}
\fancyhf{}
\rhead{Max Marks: 50}
\lhead{Exam: World War II & Decolonization}
\rfoot{Page \thepage}

\begin{document}

\begin{center}
    {\LARGE \textbf{ World War II & Decolonization }} \\[0.5em]
    \textbf{Total Maximum Marks:} 50 \quad | \quad \textbf{Total Questions:} 2
\end{center}

\vspace{1em}
\hrule
\vspace{1.5em}

\noindent \textbf{Instructions:} Answer all questions carefully. Show all mathematical derivations, intermediate steps, and calculations where applicable.

\vspace{1.5em}

\begin{enumerate}[label=\textbf{Q\arabic*.}, leftmargin=2em]

    \item \textbf{[25 Marks]} Analyse the causes and consequences of the end of World War II in 1945, focusing on the surrender of Axis powers and the establishment of the United Nations. Include evaluation of the UN Charter's principles and its impact on the post-war international order. Use primary source extracts (e.g., the Preamble of the UN Charter) to support your answer.
    
    \vspace{1em}

    \item \textbf{[25 Marks]} Examine the Indian Independence Act of 1947. Discuss the political, social, and economic factors leading to the partition of British India, and assess the immediate and long-term consequences for India and Pakistan. Incorporate analysis of contemporary perspectives (e.g., speeches by Jawaharlal Nehru and Muhammad Ali Jinnah) and evaluate the act's significance in the broader context of decolonization.
    
    \vspace{1em}

\end{enumerate}

\end{document}